\documentclass[11pt]{article}
\usepackage{labsheet_common_th}

\title{}
\date{}

\begin{document}
\thispagestyle{fancy}

\labtitle{AI Training -- บทปฏิบัติการที่ 2}{ตรวจจับข้อบกพร่องของรางรถไฟด้วย AI}{แนะนำการแบ่งส่วนภาพ (image segmentation) และการปรับจูนโมเดลแบบลงมือทำจริง}
\nocodebadge

คุณจะทำตามขั้นตอนแบบคัดลอก-วางในเทอร์มินัลและแก้ไขไฟล์ตั้งค่าแบบข้อความล้วนเพียงไฟล์เดียว
เท่านั้นเอง -- โค้ด AI ทั้งหมดมีคนเขียนไว้ให้เรียบร้อยแล้ว

\section*{สิ่งที่คุณจะได้เรียนรู้}
\begin{itemize}
  \item ความหมายของ ``Non-Destructive Testing'' (NDT หรือการทดสอบแบบไม่ทำลาย) และเหตุใด AI จึงช่วยได้
  \item ความหมายของการ ``แบ่งส่วนภาพ'' (segmentation) ด้วย AI (การระบุตำแหน่งทีละพิกเซล) แทนที่จะติดป้ายกำกับทั้งภาพ
  \item ค่าตั้งค่าเพียงค่าเดียว -- \textbf{ค่าเกณฑ์การตัดสินใจ (decision threshold)} -- ควบคุมจุดสมดุลระหว่างการตรวจจับข้อบกพร่องให้ได้ทุกจุดกับการหลีกเลี่ยงการแจ้งเตือนผิดพลาดได้อย่างไร
  \item วิธีอ่านค่าความแม่นยำ (IoU) และเปรียบเทียบกับเพื่อนร่วมชั้น
\end{itemize}

\section*{ภาพรวมของงาน}
ทีมซ่อมบำรุงทางรถไฟตรวจสอบพื้นผิวรางเพื่อหาข้อบกพร่อง (รอยแตกร้าว, การกะเทาะของผิวราง,
สนิม) ที่อาจก่อให้เกิดปัญหาหากไม่ถูกตรวจพบ -- โดยดั้งเดิมทำโดยเจ้าหน้าที่เดินตรวจตามรางแล้วดู
ด้วยตา NDT หมายถึงการตรวจสอบหาข้อบกพร่องโดยไม่ทำลายราง (เช่น แค่ถ่ายภาพเท่านั้น) ซึ่งตรงกับ
สิ่งที่ชุดข้อมูลของเราทำอยู่พอดี

\begin{notebox}
ต่างจากบทปฏิบัติการที่ 1 \textbf{ไม่มีใครต้องติดป้ายกำกับอะไรในบทนี้เลย} -- ชุดข้อมูลมาพร้อมภาพ
รางรถไฟจริง 113 ภาพอยู่แล้ว แต่ละภาพมีแผนที่ ``ค่าอ้างอิงจริง'' (ground truth) ที่วาดด้วยมือ
แสดงตำแหน่งของข้อบกพร่องอย่างชัดเจน (วาดโดยนักวิจัยผู้เผยแพร่ชุดข้อมูลนี้) งานของคุณคือฝึก AI
ให้วาดแผนที่แบบเดียวกันด้วยตัวเองบนภาพใหม่ และทำให้ AI นั้นแม่นยำที่สุดเท่าที่จะทำได้
\end{notebox}

\section*{สิ่งที่เตรียมไว้ให้คุณแล้ว}
\begin{itemize}
  \item ชุดข้อมูล (ภาพ 113 ภาพ + แผนที่ข้อบกพร่องอย่างเป็นทางการ) ถูกดาวน์โหลดไว้ใน \path|datasets/raw/rail_ndt/rsdds/| แล้ว
  \item โค้ด AI ทั้งหมด -- การสกัดฟีเจอร์, การฝึกโมเดล, การให้คะแนน, การพล็อตกราฟ -- ถูกเขียนไว้แล้วใน \texttt{scripts/} คุณมีหน้าที่แค่ \emph{รัน} เท่านั้น ไม่ต้องแก้ไข
  \item ไฟล์เดียวที่คุณจะแก้ไขคือ \texttt{configs/rail\_ndt\_config.yaml} ซึ่งเป็นไฟล์ตั้งค่าแบบข้อความล้วน (ไม่ใช่โค้ด)
\end{itemize}

\section*{ภาพรวมขั้นตอนการทำงานทั้งหมด}
\begin{center}
\resizebox{\textwidth}{!}{%
\begin{tikzpicture}[node distance=6mm]
  \node[autobox] (data) {ภาพ RSDDs + ค่าอ้างอิงจริง (ดาวน์โหลดไว้แล้ว)};
  \node[actionbox, right=of data] (config) {แก้ไข\\\texttt{rail\_ndt\_\allowbreak config.yaml}};
  \node[autobox, right=of config] (train) {ฝึกโมเดล\\\texttt{rail\_\allowbreak train.py}};
  \node[autobox, right=of train] (evaluate) {ประเมินโมเดล\\\texttt{rail\_\allowbreak evaluate.py}};
  \node[scorebox, right=of evaluate] (score) {คะแนนของคุณ:\\Mean IoU};

  \draw[flowarrow] (data) -- (config);
  \draw[flowarrow] (config) -- (train);
  \draw[flowarrow] (train) -- (evaluate);
  \draw[flowarrow] (evaluate) -- (score);

  \draw[loopB] (score.south) to[out=-55, in=-125, looseness=0.8]
    node[looplabel, midway, below=1mm] {แค่เปลี่ยนค่าเกณฑ์ -- ไม่ต้องฝึกใหม่ (Part 3)} (evaluate.south);
  \draw[loopA] (score.north) to[out=55, in=125, looseness=0.9]
    node[looplabel, midway, above=1mm] {เปลี่ยนค่าตั้งค่าโมเดล -- ต้องฝึกใหม่ (Part 4)} (config.north);
\end{tikzpicture}%
}
\end{center}
\diagramlegend

มีลูปย้อนกลับสองแบบเข้าสู่ขั้นตอนเดิม: เส้นประสีน้ำเงินใต้แผนภาพคือทางลัดที่รวดเร็ว (แค่รัน
ขั้นตอนประเมินใหม่); เส้นทึบสีส้มด้านบนต้องเริ่มใหม่ตั้งแต่การฝึกโมเดล

\section*{Part 0: ตั้งค่าเริ่มต้น (ทำครั้งเดียว)}
เปิดเทอร์มินัล (บน Mac: แอป \textbf{Terminal}; บน Windows: \textbf{Anaconda Prompt} หรือ
\textbf{PowerShell}) แล้วรันคำสั่งต่อไปนี้ทีละบรรทัด:

\begin{lstlisting}
cd AI-training
python3 -m venv .venv
source .venv/bin/activate      # Windows: .venv\Scripts\activate
pip install -r requirements.txt
\end{lstlisting}

หากคุณทำขั้นตอนนี้ไว้แล้วสำหรับบทปฏิบัติการที่ 1 ในโฟลเดอร์เดียวกัน คุณข้ามไปทำ Part 1 ได้เลย
(แค่อย่าลืมรันบรรทัด \texttt{source .venv/bin/activate} ใหม่ทุกครั้งที่เปิดหน้าต่างเทอร์มินัล
ใหม่)

\section*{Part 1: ฝึกโมเดลของคุณ}
เปิดไฟล์ \texttt{configs/rail\_ndt\_config.yaml} ด้วยโปรแกรมแก้ไขข้อความล้วนใด ๆ (Notepad,
TextEdit, VS Code -- อะไรก็ได้ที่ไม่ใช่ Word) คุณจะเห็นค่าตั้งค่าแบบนี้:

\begin{lstlisting}
model_type: random_forest
n_estimators: 100
max_depth: 8
random_seed: 42
threshold: 0.5
\end{lstlisting}

สำหรับตอนนี้ ให้คงค่าทุกอย่างไว้เหมือนเดิมแล้วรัน:

\begin{lstlisting}
python3 scripts/rail_train.py
\end{lstlisting}

คำสั่งนี้จะพิจารณาทุกพิกเซลของภาพฝึกและเรียนรู้ว่า ``พิกเซลที่เป็นข้อบกพร่อง'' มีลักษณะอย่างไร
(ความสว่าง, ลวดลายพื้นผิว, และตำแหน่งในภาพ) เทียบกับพื้นผิวรางปกติ ใช้เวลาไม่ถึงหนึ่งนาที

\section*{Part 2: ประเมินโมเดลของคุณ}
\begin{lstlisting}
python3 scripts/rail_evaluate.py
\end{lstlisting}

คำสั่งนี้จะทดสอบโมเดลของคุณกับภาพที่ไม่เคยใช้ฝึกมาก่อน แล้วแสดงผลลัพธ์ประมาณนี้:

\begin{lstlisting}
Mean IoU:  0.25 (higher is better, 1.0 = perfect)
Mean Dice: 0.38 (higher is better, 1.0 = perfect)
\end{lstlisting}

\textbf{IoU (Intersection over Union)} วัดว่าพื้นที่ข้อบกพร่องที่โมเดลของคุณทายซ้อนทับกับ
พื้นที่ข้อบกพร่องจริงมากแค่ไหน 1.0 = ซ้อนทับกันสมบูรณ์, 0 = ไม่ซ้อนทับกันเลย นี่คือคะแนนของคุณ
-- \textbf{ยิ่งสูงยิ่งดี}

นอกจากนี้ยังบันทึกภาพสองภาพไว้ด้วย:
\begin{itemize}
  \item \texttt{results/figures/rail\_eval.png} -- ภาพทดสอบหลายภาพเรียงกัน แต่ละภาพแสดง
  พื้นที่ข้อบกพร่องจริงและที่ทายซ้อนทับลงบนภาพรางรถไฟจริงด้วยสี -- \textbf{สีเหลือง} =
  ตรวจจับถูกต้อง, \textbf{สีเขียว} = ข้อบกพร่องจริงที่โมเดลพลาดไป, \textbf{สีแดง} = การแจ้งเตือน
  ผิดพลาด -- พร้อมค่า IoU ของภาพนั้นกำกับไว้ด้านบน ลองดูภาพนี้ -- เป็นวิธีที่มีประโยชน์ที่สุดในการ
  ทำความเข้าใจว่าโมเดลของคุณกำลังทำอะไรอยู่
  \item \texttt{results/figures/rail\_eval\_threshold\_curve.png} -- ค่า Mean IoU และ Mean
  Dice ที่คำนวณอัตโนมัติทุกค่าเกณฑ์ตั้งแต่ 0.05 ถึง 0.95 สำหรับโมเดลที่ฝึกไว้ตัวเดียวกันนี้
  พร้อมทำเครื่องหมายค่าเกณฑ์ปัจจุบันของคุณไว้ เป็นตัวอย่างล่วงหน้าของจุดสมดุลที่คุณจะลองปรับ
  ด้วยตัวเองใน Part 3
\end{itemize}

\section*{Part 3: ปรับค่าเกณฑ์ (เครื่องมือหลักที่ ``ไม่ต้องฝึกโมเดลใหม่'')}
โมเดลของคุณไม่ได้แค่บอกว่า ``เป็นข้อบกพร่อง'' หรือ ``ไม่ใช่'' สำหรับแต่ละพิกเซล แต่ให้คะแนน
ความมั่นใจตั้งแต่ 0 ถึง 1 ค่าตั้งค่า \texttt{threshold} (ค่าเกณฑ์) เป็นตัวกำหนดว่าโมเดลต้อง
มั่นใจแค่ไหนก่อนจะทำเครื่องหมายพิกเซลนั้นว่าเป็นข้อบกพร่อง:

\begin{center}
\begin{tabular}{@{}lp{9.5cm}@{}}
\toprule
\textbf{threshold (ค่าเกณฑ์)} & \textbf{ผลลัพธ์} \\
\midrule
ต่ำ (เช่น 0.2)  & ทำเครื่องหมายพิกเซลมากขึ้น -- จับข้อบกพร่องจริงได้มากขึ้น แต่ก็แจ้งเตือนผิดพลาดมากขึ้นด้วย \\
สูง (เช่น 0.8) & ทำเครื่องหมายพิกเซลน้อยลง -- แจ้งเตือนผิดพลาดน้อยลง แต่เสี่ยงพลาดข้อบกพร่องจริง \\
\bottomrule
\end{tabular}
\end{center}

เปิด \texttt{configs/rail\_ndt\_config.yaml} เปลี่ยนเฉพาะค่า \texttt{threshold} แล้วรัน
\textbf{เฉพาะ}ขั้นตอนประเมินใหม่ (ไม่ต้องฝึกโมเดลใหม่):

\begin{lstlisting}
python3 scripts/rail_evaluate.py
\end{lstlisting}

\begin{tipbox}
ลองหลายค่า (เช่น 0.3, 0.4, 0.5, 0.6, 0.7, 0.8) แล้วจดค่า Mean IoU ที่ได้แต่ละครั้ง นี่คือ
วิธีที่เร็วที่สุดในการปรับปรุงคะแนนของคุณ และไม่ต้องฝึกโมเดลใหม่ -- แต่ละครั้งใช้เวลาเพียงไม่กี่
วินาที
\end{tipbox}

\section*{Part 4: (ทางเลือกเพิ่มเติม/ขั้นสูง) ฝึกโมเดลใหม่ด้วยค่าตั้งค่าอื่น}
หากต้องการลองเพิ่มเติม ให้แก้ไข \texttt{model\_type}, \texttt{n\_estimators},
\texttt{max\_depth} หรือ \texttt{random\_seed} แล้วรันทั้งสองขั้นตอนใหม่ตามลำดับ (ค่าตั้งค่า
เหล่านี้จะมีผลก็ต่อเมื่อฝึกโมเดลใหม่เท่านั้น):

\begin{lstlisting}
python3 scripts/rail_train.py
python3 scripts/rail_evaluate.py
\end{lstlisting}

จากนั้นทำ Part 3 ซ้ำเพื่อปรับค่าเกณฑ์ใหม่สำหรับโมเดลใหม่ของคุณ -- ค่าเกณฑ์ที่ดีที่สุดอาจ
แตกต่างกันไปตามค่าตั้งค่าที่ต่างกัน

\section*{แข่งกับเพื่อนร่วมชั้น}
จดบันทึกคะแนนไว้:

\begin{center}
\small
\begin{tabular}{@{}*{5}{>{\raggedright\arraybackslash}p{2.9cm}}@{}}
\toprule
\textbf{ชื่อของคุณ} & \textbf{model\_type} & \textbf{ค่าตั้งค่าหลัก} & \textbf{threshold} & \textbf{Mean IoU (ยิ่งสูงยิ่งดี)} \\
\midrule
 & & & & \\[10pt]
 & & & & \\[10pt]
 & & & & \\[10pt]
\bottomrule
\end{tabular}
\end{center}

\section*{Part 5 (ทางเลือกเพิ่มเติม): สร้างรายงานที่แชร์ได้}
เมื่อคุณพอใจกับผลลัพธ์ที่ได้แล้ว ให้แปลงเป็นไฟล์ PDF หน้าเดียวที่พิมพ์หรือส่งให้เพื่อนร่วมชั้น/
วิทยากรได้:
\begin{lstlisting}
python3 scripts/rail_report.py
\end{lstlisting}
คำสั่งนี้จะอ่านผลลัพธ์ล่าสุดของคุณแล้วบันทึกเป็น
\texttt{results/reports/rail\_report.pdf} -- รวมค่าตั้งค่า คะแนน ตัวอย่างการทาย และกราฟ
ค่าเกณฑ์ไว้ในหน้าเดียว รันซ้ำได้ทุกเมื่อหลัง Part 2 (หรือหลังปรับค่าเกณฑ์ใน Part 3) เพื่อบันทึก
ผลลัพธ์ที่ดีที่สุดของคุณในปัจจุบัน

\section*{อภิธานศัพท์}
\begin{itemize}
  \item \textbf{NDT (Non-Destructive Testing)} -- การตรวจสอบหาข้อบกพร่องโดยไม่ทำให้สิ่งของนั้นเสียหาย เช่น ตรวจจากภาพถ่ายแทนที่จะตัดรางออกมาดู
  \item \textbf{Segmentation (การแบ่งส่วนภาพ)} -- แทนที่จะติดป้ายกำกับทั้งภาพว่า ``มีข้อบกพร่อง'' หรือไม่ โมเดลจะทำเครื่องหมาย \emph{พิกเซลใดบ้าง} ที่เป็นส่วนหนึ่งของข้อบกพร่อง
  \item \textbf{Ground truth (ค่าอ้างอิงจริง)} -- คำตอบที่ถูกต้อง ในที่นี้คือแผนที่ที่วาดด้วยมือแสดงตำแหน่งของข้อบกพร่องแต่ละจุดอย่างแม่นยำ จัดทำโดยนักวิจัยเจ้าของชุดข้อมูลต้นฉบับ
  \item \textbf{Threshold (ค่าเกณฑ์)} -- จุดตัดความมั่นใจ (0 ถึง 1) ที่เมื่อเกินค่านี้พิกเซลจะถูกเรียกว่า ``ข้อบกพร่อง'' เป็นค่าตั้งค่าเดียวที่คุณเปลี่ยนได้โดยไม่ต้องฝึกโมเดลใหม่
  \item \textbf{IoU (Intersection over Union)} -- พื้นที่ซ้อนทับระหว่างพื้นที่ข้อบกพร่องที่โมเดลของคุณทายกับพื้นที่จริง หารด้วยพื้นที่รวมของทั้งสอง 1.0 = เหมือนกันทุกประการ, 0 = ไม่ซ้อนทับกันเลย
  \item \textbf{Dice score} -- ค่าวัดการซ้อนทับที่คล้ายกับ IoU มาก มักมีค่าสูงกว่าเล็กน้อยสำหรับการทายแบบเดียวกัน ทั้งสองค่าถูกรายงานไว้เพื่อให้คุณเปรียบเทียบได้
  \item \textbf{False positive (ผลบวกลวง)} -- พิกเซลปกติที่ถูกทำเครื่องหมายผิดว่าเป็นข้อบกพร่อง (การแจ้งเตือนผิดพลาด)
  \item \textbf{False negative (ผลลบลวง)} -- พิกเซลข้อบกพร่องจริงที่โมเดลพลาดไป
  \item \textbf{Class imbalance (ความไม่สมดุลของคลาส)} -- ในทุกภาพ พิกเซลส่วนใหญ่คือรางปกติ มีเพียงพื้นที่เล็ก ๆ เท่านั้นที่เป็นข้อบกพร่อง โมเดลถูกออกแบบมาโดยเฉพาะเพื่อรองรับเรื่องนี้ มิฉะนั้นโมเดลอาจได้คะแนนสูงอย่างหลอกลวงเพียงเพราะไม่เคยทาย ``ข้อบกพร่อง'' เลย
  \item \textbf{Random Forest} -- โมเดลที่ประกอบด้วย ``ต้นไม้ตัดสินใจ'' แบบง่าย ๆ จำนวนมากที่ออกเสียงโหวตแต่ละพิกเซล \texttt{n\_estimators} = จำนวนต้นไม้, \texttt{max\_depth} = ความละเอียดของคำถามที่แต่ละต้นถามได้
\end{itemize}

\begin{warnbox}
\begin{itemize}
  \item \textbf{``no RSDDs images found''} -- รัน \texttt{python3 scripts/fetch\_rsdds.py} ก่อน เพื่อดาวน์โหลดชุดข้อมูล (ทำครั้งเดียวก็พอ)
  \item \textbf{``no trained model found''} -- รัน \texttt{python3 scripts/rail\_train.py} ก่อน \texttt{rail\_evaluate.py}
  \item \textbf{เปลี่ยนค่าตั้งค่าแล้วแต่คะแนนไม่เปลี่ยน} -- มีเพียง \texttt{threshold} เท่านั้นที่มีผลทันที ค่าตั้งค่าอื่นทั้งหมดต้องรัน \texttt{rail\_train.py} ใหม่ก่อน
  \item \textbf{คะแนนแย่กว่าของเพื่อนร่วมชั้นมาก ทั้งที่ใช้ค่าตั้งค่าเหมือนกันทุกประการ} -- ตรวจสอบว่า \texttt{random\_seed} เหมือนกันหรือไม่ เพราะค่านี้กำหนดว่าภาพใดจะอยู่ในกลุ่มฝึกหรือกลุ่มทดสอบ ค่า seed ที่ต่างกันหมายถึงการฝึก/ทดสอบด้วยชุดภาพที่แบ่งต่างกันไปเลย
\end{itemize}
\end{warnbox}

\section*{ประเด็นให้ลองคิดต่อ (พูดคุยกับเพื่อนร่วมชั้น)}
\begin{itemize}
  \item เหตุใดการทำคะแนน IoU สูงจึงยากกว่าการทำคะแนนความแม่นยำ (accuracy) สูงในกรณีนี้ เมื่อพิจารณาว่าพื้นที่ข้อบกพร่องมีขนาดเล็กมากเมื่อเทียบกับภาพทั้งภาพ?
  \item ในระบบตรวจสอบจริง คุณอยากให้ใช้ค่าเกณฑ์ต่ำ (แจ้งเตือนผิดพลาดมากขึ้น) หรือค่าเกณฑ์สูง (พลาดข้อบกพร่องมากขึ้น)? คำตอบของคุณเปลี่ยนไปหรือไม่ หากรอยแตกร้าวที่พลาดไปอาจทำให้รถไฟตกราง?
  \item ชุดข้อมูลนี้ไม่มีภาพ ``รางสะอาด'' ที่ไม่มีข้อบกพร่องเลยแม้แต่ภาพเดียว -- ทุกภาพมีข้อบกพร่องอย่างน้อยหนึ่งจุด สิ่งนี้อาจส่งผลต่อสิ่งที่โมเดลเรียนรู้ไปอย่างไร และคุณอยากเพิ่มอะไรเพื่อแก้ไขปัญหานี้?
\end{itemize}

\end{document}
